\chapter{SMC Software architecture}
\section{Robot classes}
Ideally, we want the same controller to run on different robots.
This is achieved by having generic robot classes,
and using generic calls to get things like Jacobians.
Of course, different types of robots have different additional
features and functions, but as much as possible is shared.
For example mobile robots have an ``end-effector'',
which is their base position, and this allows for 
use of the same Cartesian controllers on the mobile robot
as on manipulators (otherwise you'd have to have 2 identical functions
for path following for example).

Importantly, if applicable, a robot can ``switch'' between 
these classes by selecting a \fbox{\texttt{control\_mode}}.
This means that if one sets \fbox{\texttt{base\_only}} mode
on a mobile robot with an arm, the robot will only give
its base's position, velocity and Jacobian. If a Cartesian controller
is called, the goal error will be computed for the base's position
instead of the end-effector, and the actuation will be commanded
only to the base instead of the whole robot (all other control inputs
are set to 0).
Of course, some controllers don't make sense unless a particular
kind of robot and control mode are used - in the code this is denoted
with type assertions.

In code, the classes of robots are referred to as interfaces.
This is because capabilities are, to a certain extend, modular. For example
a dual arm mobile robot has the same mobile base interface as a single arm mobile robot,
and this structure allows for less code.
The existing interfaces robots are:
\begin{itemize}
		\item \fbox{\texttt{MobileBaseInterface}} - functions for a base planar joint 
		\item \fbox{\texttt{SingleArmInterface}} - functions for a kinematic chain with a single end-effector
		\item \fbox{\texttt{ForceTorqueSensorInterface}} - functions for managing a force/torque sensor like 
				basic low-pass filtering. Combined with SingleArmInterface 
				to form ForceTorqueOnSingleArmWrist for manipulators
				with a force/torque sensor embedded in the wrist of the arm
				(has functions to map the wrench to the end-effector frame etc)
		\item \fbox{\texttt{DualArmInterface}} - functions for managing a dual arm robot. Inherits 
				from SingleArmInterface as that allows controllers for SingleArmInterface
				to be applied to the absolute frame via the absolute-relative
				frame transformations. Additionally, each arm can be controlled individually,
				or both simultaneously
				via the corresponding control modes.
		\item \fbox{\texttt{WholeBodySingleArmInterface}} - functions for a mobile robot with a single manipulator.
				The base and the arm can be controlled separately or together.
		\item \fbox{\texttt{WholeBodyDualArmInterface}} - functions for a mobile robot with a single manipulator.
				The base and (either of) the arms can be controlled separately or together.
\end{itemize}

\subsection{AbstractRobotManager}
AbstractRobotManager is the base class for all robots. 
It provides the most generic functionality,
and specifies what specific robots need to have implemented.
The table of relevant generic functions 
is given in \ref{table-abstract-robot} (code bookkeeping like
starting the visualizer or connecting to the robot is skipped).

\begin{table}[h!]
\centering
\label{table-abstract-robot}
\begin{tabular}{c | p{9cm}}
\textbf{Attribute/Function} & \textbf{Description}\\
\hline 
control\_mode & enum for determining control mode. 
options are whole\_body, base\_only, upper\_body, 
	   left\_arm\_only, right\_arm\_only \\
\hline 
model & Pinocchio model of the robot. this is a getter function via Python's @property  \\
\hline 
data & Pinocchio data for the robot. this is a getter function via Python's @property  \\
\hline 
\_args & all arguments, including defaults. works
as a Config struct\\
\hline 
q & vector of joint positions, $ q \in \mathbb{R}^{nq}  $. this is a getter function via Python's @property \\
\hline 
v & vector of joint velocities, $ v \in \mathbb{R}^{nv}  $. this is a getter function via Python's @property  \\
\hline 
comfy\_configuration & decent hand-defined home position \\
\hline 
forwardKinematics() & updates frame positions. called automatically in ControlLoopManager  \\
\hline 
getJacobian() & returns the Jacobian. Always in the body frame. The actual frame depends on the control mode.   \\
\hline 
sendVelocityCommand(v) & sends command to real robot, or integrates by 
$ \Delta t  $ in simulation. called automatically in ControlLoopManager  \\
\hline 
stopRobot() & stops the robot however it should (sometimes sending $ v =0  $ is not enough) \\
\hline 
setFreedrive() and stopFreedrive() & toggle freedrive, or leadthrough - whatever allows you to move
the robot manually \\
\hline 
openGripper() and closeGripper() & open/close gripper. usually build upon setPosition(p) 
in the gripper driver, but in nearly all cases we just need this\\
\end{tabular}
\caption{Generic functions available on all robots}
\end{table}

\subsection{Mobile base}
There's a differential drive base and an omnidirection base.
The differential drive is modelled as an $ SE (2)  $ joint with
linear velocity control input for the $ y  $ direction (pointing forward)
and rotational velocity control input $ \dot{\theta}  $, i.e.
\begin{equation}
u_{ \text{diff drive} } = 
\begin{bmatrix} 
v_{ x } \\
0\\
\dot{\theta}
\end{bmatrix} 
\end{equation}
with position being 
$ q = \begin{bmatrix} x & y & \theta \end{bmatrix}^{ T }   $. 
Importantly, Pinocchio treats instead uses a complex number for
rotations in 
$ SE (2)  $, i.e. 
$ q = \begin{bmatrix} x & y & \cos(\theta) & \sin(\theta) \end{bmatrix}^{ T }      $
%The reason we model the differential drive this way is because
%that's the control interface we get on the robot Heron in the lab
%(we don't get wheel access). But 
By assuming the no-slip setting, controlling
wheel speeds or the desired twist are equivalent.

The fully actuated base (ex. pseudo-omnidirectional) uses
\begin{equation}
u_{ \text{fully actuated} } = 
\begin{bmatrix} v_{ x }\\ v_{ y } \\ \dot{\theta} \end{bmatrix} 
\end{equation}

\begin{table}[h!]
\centering
\label{table-mobile-base-interface}
\begin{tabular}{c | p{9cm}}
\textbf{Attribute/Function} & \textbf{Description}\\
\hline 
base\_SE2\_pose() & returns position as 
$ \begin{bmatrix} x & y & \theta \end{bmatrix}   $\\
\hline 
T\_w\_b & returns the transformation from the world
frame $ \left\{ w \right\}   $ to the end-effector frame 
$ \left\{ b \right\}   $, i.e. $ T_{ w }^{ b } \in SE (3) $. it's a getter via Python's @property\\
\hline 
T\_w\_e & returns the same as T\_w\_b, but represents an ``end-effector'' frame.
it's an evil hack to use some SingleArmInterface controllers with no additional code\\
\hline 
getJacobian() & returns 
$   J = \begin{bmatrix} 
		R_{ w }^{ b } & 0 \\
		0 & 1
\end{bmatrix} $,
meaning it maps velocities in the base's frame $ \left\{ b \right\}   $ 
to world velocities.
\\
\end{tabular}
\caption{Mobile robot specific functions.}
\end{table}

\paragraph{How to run}
Run an applicable script with the following arguments:
\begin{itemize}
		\item just the diff drive base: \fbox{\texttt{\ttfamily--robot mir}}
		\item just the fully actuated base: \fbox{\texttt{--robot mir  --fully-actuated-base}} 
		\item diff-drive on single arm mobile robot: \fbox{\texttt{--robot mir --mode base\_only}}
		\item fully-actuated on dual arm mobile robot: \fbox{\texttt{--robot myumi --mode base\_only}}
\end{itemize}


\subsection{Single arm interface}
The SingleArmInterface covers the functionality of single arm robots (manipulators), 
i.e. of articulated robots. 
None of the controllers assume 6 dofs,
and there are controllers that exploit the additional degrees of freedom
which will be covered in the controllers section.

The relevant functions are given in table \ref{table-single-arm-interface}.

\begin{table}[h!]
		\centering
		\label{table-single-arm-interface}
		\begin{tabular}{c|p{7cm}}
\textbf{Attribute/Function} & \textbf{Description}\\
\hline 
T\_w\_e & returns the transformation from the world
frame $ \left\{ w \right\}   $ to the end-effector frame 
$ \left\{ e \right\}   $, i.e.$ T_{ w }^{ e } \in SE (3) $. it's a getter via Python's @property\\
\hline 
getJacobian() & returns the \textbf{body Jacobian} $ J \in \mathbb{R}^{6 \times nv}  $
from the world to the end-effector frame \\
\hline 
computeManipulabilityIndex() & computes the manipulability index
$ m = \sqrt{\det (J J^{ T })}  $ \\
\hline 
computeManipulabilityIndexQDerivative() & computes the derivative of the manipulability index
$ \frac{\partial m}{\partial q}   $ via the derivatives of forward kinematics\\
		\end{tabular}
		\caption{Functions of the single arm interface}
\end{table}


\paragraph{How to run}
Run an applicable script with the following arguments:
\begin{itemize}
		\item for UR5e robot \fbox{\texttt{--robot ur5e}}
		\item for one of YuMi's arms \fbox{\texttt{--robot myumi --control-mode left\_arm\_only}}
\end{itemize}

\subsection{Dual arm robots}
The DualArmInterface covers the functionality of dual arm robots,
which are simply 2 manipulators.
None of the controllers assume 6 dofs per arm,
and there are controllers that exploit the additional degrees of freedom
which will be covered in the controllers section.
The left arm is referred to with the index $ l  $, and the right with $ r  $.
In the code we assume that both arms are the same and have
the same number of dofs $ nql = nqr  $.

Each arm can be controlled separately via their control modes,
\fbox{\texttt{left\_arm\_only}} and \fbox{\texttt{right\_arm\_only}}.
This means that \fbox{\texttt{q}} returns only the joint configuration of the left or the right arm,
\fbox{\texttt{T\_w\_e}} returns $ T_{ w }^{ l }  $ or $ T_{ w }^{ r }  $ respectively and so on.
Attributes like \fbox{\texttt{nq}} and \fbox{\texttt{nv}} function accordingly.

The two arms can be controlled simultaneously with
control mode \fbox{\texttt{upper\_body}},
and if the robot isn't mobile also with \fbox{\texttt{whole\_body}}.
In this case $ q  $ returns the joint configuration of both arms concatenated together.
The convention is that the left arm joints are given first,
i.e. 
$ q = \begin{bmatrix} q_{ l_{ 1 } } & \dots q_{ l_{ nq l } } & 
q_{ r_{ 1 } } & \dots  & q_{ nqr }\end{bmatrix} ^{ T }   $.
The Jacobian is combined arm Jacobian and is also simply the concatenation of
the single arm Jacobians:
\begin{equation}
\begin{bmatrix} 
v_{ l } \\
v_{ r }
\end{bmatrix} =
\underbrace{\begin{bmatrix} J_{ l } & 0 \\
		0 & J_{r }
\end{bmatrix}}_{ J } 
\begin{bmatrix} 
q_{ l } \\
q_{ r }
\end{bmatrix} 
\end{equation}

However, the most useful way is to control the robot with the concept of an absolute frame
$ T_{ w }^{ abs }  $:
\begin{equation}
T_{ w }^{ abs } = \text{interpolate } (T_{ w }^{ l }, T_w^{ r }, 0.5)
\end{equation}
where interpolate() interpolates between the first and the second
$ SE (3)  $ object with the parameter $ s \in [0,1]  $. 
In other words, the absolute frame is ``middle frame'' between the two arms.
Often, we want to specify the end-effector poses of the arms in relation
to the absolute frame. The transformation between the absolute frame and the
end-effector frames is called the relative frame. It does not need to be symmetric.
In the code we use $ T_{ abs }^{ l }  $ and $ T_{ abs }^{ r }  $. If these two frames
are fixed, we can use a single arm controller strategy on the absolute frame,
and let the code map to both arms via their relative transformations behind the scenes.
In the cart-pulling problem, this means we can define the desired trajectory for
the handlebar as the desired trajectory of the absolute frame,
and automatically get the trajectories for the left and right end-effector 
since both grasp the handlebar rigidly.
More on this in the control section.


The relevant functions are given in table \ref{table-dual-arm-interface}.

\begin{table}[h!]
		\centering
		\label{table-dual-arm-interface}
		\begin{tabular}{c|p{7cm}}
\textbf{Attribute/Function} & \textbf{Description}\\
\hline 
T\_w\_e & depending on the mode, it returns
$ T_{ w }^{ l }, T_{ w }^{ r }  $ or $ T_{ w }^{ abs }  $. 
it's a getter via Python's @property\\
\hline 
[TODO] random convenience function & random convenience function description
		\end{tabular}
		\caption{Functions of the dual arm interface}
\end{table}

\subsection{Single arm mobile robots}
These are manipulators on a mobile base. They are referred to as
WholeBodySingleArmInterface.
All functions work as expected depending on the control mode.
\paragraph{[TODO]: write Jacobians for each control mode (already in code, but needs to be written)}

\subsection{Dual arm mobile robots}
These are dual manipulators on a mobile base. They are referred to as
WholeBodyDualArmInterface.
All functions work as expected depending on the control mode.
\paragraph{[TODO]: write Jacobians for each control mode (already in code, but needs to be written)}

\subsection{ForceTorqueSensorInterface}
TODO: writeabout it
\begin{itemize}
		\item bias estimation on by default, happens every time you start 
		\item f/t sensor filtering is moving average AND THEN a low-pass filter
		\item atm we only have direct force feedback control
\end{itemize}


\section{Controllers}
\subsection{ControlLoopManager}
Every controller is written as a function which is inserted into 
\fbox{\texttt{ControlLoopManager}}, which handles real-time-plotting,
logging, communication with the robot if running on hardware,
and ensures the commands are sent at a prescribed frequency ---
it simply calls \fbox{\texttt{sleep(1/freq - time\_spent\_computing)}}.
In other words, controllers are managed with composition over inheritance.

Every command velocity is saturated (clipped) to the maximum velocity allowed by the robot.
This is defined in RobotManager.sendVelocityCommand() and is always applied.

An important feature of ControlLoopManager is \fbox{\texttt{history}} or
\fbox{\texttt{PastWindow}} which is a dictionary (hash map) of circular queues keeping 
whatever quantities the user selects. For example,
in the cart-pulling case, we want to keep track of the path the base traversed.
The size of the past window is determined by the 
\fbox{\texttt{--past-window-size}} argument. It determines the number
of iterations saved, meaning that the actual time saved depends on the control
frequency as well.

Logging is handled similarly, with a dictionary of lists (linked lists)
which save whatever quantity the user selects.
These can be pickled (Python's serialization module) to a file.
Note that the whole of Logger class is pickled so that metadata
is also stored conveniently.
There is code to plot logs, and it assumes that what's saved are vectors of 
floating points. 
In total, it looks like the following:

\begin{algorithm}[H]
\caption{ControlLoopManager.run()}
\begin{algorithmic}[1] % The [1] adds line numbers
\Procedure{ControlLoopManager.run}{cfg, robot\_manager, control\_loop}
\While{not breakFlag}
    \State start\_time $\gets$ current\_time()
    \State robot\_manager.\_step() - \textit{updates kinematic chain etc}
    \State breakFlag, log\_item, past\_item = control\_loop(..., cfg, robot\_manager, past\_window, t)
    \State log.update(log\_item)
    \State past\_window.update(past\_item)
    \State visualizer.update() - \textit{executed in parallel process}
    \State end\_time $\gets$ current\_time()
    \State sleep(1/freq - (end\_time - start\_time))
\EndWhile
\EndProcedure
\end{algorithmic}
\label{cart-reference-algorithm}
\end{algorithm}
Importantly, ControlLoopManager listens and handles the \fbox{\texttt{CTRL+C}} (SIGINT) signal
and stops the robot before exiting.


\subsection{ControlLoopTemplates}
To avoid code copy-pasting, there is one more optional layer of composition.
The two most important templates are \fbox{\texttt{PointToPointTemplate}}
and \fbox{\texttt{PathFollowingTemplate}}. These compute 
errors to the goal point, or to the next path point respectively.
They also save the expected logging items (ex. goal error over time)
and likewise manage the history dictionary. Both can still be expanded
in the same manner.

Since this is optional, it won't be explored further here.

\subsection{GenericControlLoop} 
This serves simply to demonstrate what every control loop must have to 
fit the framework:
\begin{minted}[mathescape, linenos]{python}
def genericControlLoop(
    X: Any,
    args: Namespace,
    robot: AbstractRobotManager,
    t: int,  # will be float eventually
    past_data: dict[str, deque[np.ndarray]],
) -> tuple[bool, dict[str, np.ndarray], dict[str, np.ndarray]]:
    breakFlag = False
    log_item = {}
    save_past_item = {}

    # compute breakFlag and v_cmd however you want

    robot.sendVelocityCommand(v_cmd)

    log_item["q"] = robot.q
    return breakFlag, save_past_item, log_item
\end{minted}

To make control\_loops convenient to use when writing main files,
the idea is to write a small wrapper function which instantiates
a \fbox{\texttt{ControlLoopManager}} for the given
\fbox{\texttt{control\_loop}}.
This usually looks like
\begin{minted}[mathescape, linenos]{python}
def genericController(
    X: Any,
    args: Namespace,
    robot: AbstractRobotManager,
	run=True
) -> ControlLoopManager | None:
    """
    GenericController
	-----------------
	description
	"""
    assert type(X) == what_it_should_be
	controlLoop = functools.partial(genericControlLoop, X, args, robot)
    log_item = {
        "qs": np.zeros(robot.nq),
        "dqs": np.zeros(robot.nv),
        "dqs_cmd": np.zeros(robot.nv),
        "err_norm": np.zeros((1,)),
    }
	save_past_item = {}
    loop_manager = ControlLoopManager(robot, controlLoop, args, save_past_item, log_item)   
    if run:
         loop_manager.run()
    else:
         return loop_manager
\end{minted}
The \fbox{\texttt{partial(X,...)}} allows passing additional arguments
to \fbox{\texttt{controlLoop}} while preserving 
\fbox{\texttt{ControlLoopManager}} API.
The if-run serves as a hot-fix for interfacing with ROS2. 
More on that in Appendix \ref{sec-ros2}.


\section{Multiprocessing and networking in SMC}
Any robotics software has to be capable of parallel computation.
Simply put, you don't want your controller to stall because the processor
is busy running a visualization, or waiting on a camera update.
Normally, this would be handled with multithreading,
but since Python still has the Global Interpreter Lock (GIL),
all Python threads run on a single processor core.

For this reason, we are forced to implement parallelism via processes.
Fortunately, since SMC is a library for quick-and-dirty implementation
for testing and throwing away robotics experiments, we can get away with
spawning all required processes in the beginning, and rest assured
that they will run in the 5 minutes our experiments last
(if you want serious code for serious applications, why are you looking at a Python library?).

Interprocess communication (IPC) in SMC is designed to be as lean and as quick as possible.
In other words, \textbf{we always aim to minimize the amount of data transferred
between processes, and to always have the last piece of data available},
irrespective of the actual protocol used.

There are 3 IPC techniques in SMC: 
\begin{itemize}
\item \textbf{queues} - for asynchronous data transfer for processes
		running at substantially different frequencies,
		ex. for visualization
\item \textbf{shared memory} - for instantaneous data access
		for processes running at comparable frequencies in the same machine
\item \textbf{sockets} - for processes running at different machines,
		with serialization from protobuf (allows for heterogeneous software stacks)
\end{itemize}

From the user point of view, IPC techniques are hidden
behind basic multi-process interaction paradigms, which
are implemented in a dedicated \fbox{\texttt{ProcessManager}}.
The main process (the one spawned when running an example script) is reserved
for control. Other processes are spawned by passing a function
to be run in a separate process
to a particular \fbox{\texttt{ProcessManager}}, which will be introduced next.

\subsection{ProcessManager}
\fbox{\texttt{ProcessManager}} is a wrapper class around the function
you want to run in a separate process.
That function is referred to as \fbox{\texttt{side\_function}}.
It's sole purpose is to spawn the process and the selected IPC method
for you. It's \textbf{your job} to declare which data you're exchanging
to actually send and receive said data, in the same manner as for logging
and plotting.

The following multi-process interaction paradigms are supported in SMC:
\begin{itemize}
\item \fbox{\texttt{ConsumerProcess}} - the side function only takes
	in data from the control process via queue (ex. visualizer, plotter)
\item \fbox{\texttt{ProducerProcess}} - the side function only provides 
	data to the control process via queue (ex. camera outputs)
\item \fbox{\texttt{ProducerConsumerProcess}} - combines
	functions of the consumer and producer processes
\item \fbox{\texttt{CollaborativeProcess}} - like producer-consumer
	process, but the communication is implemented via shared memory.
	Allows for nearly-instantaneous communication,
	ex. required in path planning (provide current position,
	obtain current path)
\end{itemize}

The producer processes require an \fbox{\texttt{init\_value}} argument,
which will be return before the first producer output.
The consumer processes require an \fbox{\texttt{init\_command}}
argument to begin computation at process spawn time, instead
of when the controller starts.
Additional arguments to the side function are to be provided 
via \fbox{\texttt{partial}}-ing.

\paragraph{Using ProcessManagers}
The process with your function will be started 
immediately upon construction of a ProcessManager object. 
You then interact with the process with the following functions called
from the ProcessManager object:
\begin{itemize}
\item \fbox{\texttt{sendCommand(command)}} - sends the command/message
		to a consumer process 
\item \fbox{\texttt{getData()}} - obtain the latest information
	from the producer process - \textbf{always}
		returns the latest received data
\end{itemize}
Depending on the underlying IPC used, the side functions will need
to receive different IPC objects, which then determine
how \fbox{\texttt{sendCommand }} and \fbox{\texttt{getData}} work.
The simplest way to navigate this is to look at existing examples
and replicating the procedure therein.

\subsection{Multiprocessing examples}
\paragraph{Producer-consumer examples}
Producer-consumer is multiprocessing lingo for a communication
relationship where one process ``produces'' information (ex. messages, tasks,
commands) and the other process ``consumes'' that information.
In other words, the communication goes one way only.
Another popular name for this communication relationship is publisher-subscriber.

A hands-on example can be found in
\fbox{\texttt{examples/vision/camera\_no\_lag.py}}
In this case control is the consumer, and the side process is the producer.
The side process obtains images from a camera, processes them,
and puts the result in a queue. The results are then made available
in the control loop. This example serves only to provide
the structure of this would work, but there is no actual task.
An example of how this setup is used 
for visual servoing can be found in
\fbox{\texttt{examples/planar\_pushing\_via\_visual\_servoing}.}

This communication relationship can also be found throughout the codebase.
The control loop is a producer for both the plotter
and the visualizer. Both use queues. 
In path following control from a planner, the path planner consumes
the current position of the robot, and produces the path to be followed.
Since this needs to be as close to instantaneous as possible,
shared memory is used for both pieces of information.


\subsection{Networking}
Working with networking requires a basic operational knowledge of computer
networks (what an IP address is etc). This is unavoidable
because it is very likely you'll have to set up and debug the network yourself.
You are referred to 
\href{https://www.youtube.com/watch?v=0Hhaf4q3eqI}{an introductory video} 
if you have no experience with networking.

More info on networking in SMC can be found in 
\newline
\fbox{\texttt{smc/multiprocessing/networking/notes.md}}.
Examples of networking in SMC can be found in 
\newline
\fbox{\texttt{examples/networked\_clik\_client}}
One example covers connecting 2 SMC clients, and the other connecting 
an SMC client with a ROS1 client.




\section{Multiprocessing in ROS2}
If you are a ROS2 user, just use ROS2 for multiprocessing.
The bare-bones SMC implementation serves primarily to avoid 
ROS2 as a dependency (as this would be the only thing in it that's used).
RCL automatically assigns different IPC protocols to
its publisher-subscriber structures, and you do not have to think
about what is ``under the hood.''
If you have to use ROS2 for interfacing with your robot,
you will have to learn how it works to access all features of your robot.
In this case you are referred to \href{https://docs.ros.org/en/rolling/index.html}{ROS2 documentation},
and are advised to use only control from SMC.
This requires a specific pre-set argument list, and a particular node structure
to work. More on this in Appendix \ref{sec-ros2}.




\section{Quality of life features}
\subsection{Visualizer}
It's Meshcat. It visualizes the robot and the end-effector frame. 
It visualizes trajectories and plans as a filtered series of frames in case of poses,
or points in case of positions.
The visualizer runs as a separate process, and does not slow down control on the main process.

For the purposes of validating plans, there is a premade map with boxes as obstacles,
visualized as such. Visualizing arbitrary boxes and sphere is straight-forward,
but no convenience code is provided.

\paragraph{How to run}
Opens in the browser if the \fbox{\texttt{--visualizer}} argument is provided (on by default,
can be turned off with \fbox{\texttt{--no-visualizer}}). 

\subsection{Plotter}
There is a basic real-time plotter implemented in Matplotlib. It simply
plots the most recent logs (stored in a circular queue) in real time. 
The plotter runs as a separate process, and does not slow down control on the main process.

\paragraph{How to run}
Opens in a new window if the \fbox{\texttt{--plotter}} argument is provided (on by default,
can be turned off with \fbox{\texttt{--no-plotter}}). 
Each new control loop spawns a new plotter (destroying the previous window
and then opening a new one), which can be a bit annoying.





\section{Suggested workflow}
\label{sec-suggested-workflow}
Here a typical workflow to developing and testing a novel program,
possibly with a new controller, is outlined.
To begin, a typical main file is described.

\subsection{Writing a new program/script}
\subsubsection{Main file structure}
\begin{minted}[mathescape, linenos]{python}
from smc import load_config, getRobotFromConfig
from smc.control.some_module.some_controller import someController

if __name__ == "__main__":
    cfg = load_config()
    robot = getRobotFromConfig(cfg)

    #  your mainfile code goes here
    somecontroller(cfg, robot, ...)

    # set expected behaviour here:
    # the library can't know where the end of a mainfile is - 
    # so you have to have this here
    if cfg.real:
        robot.stopRobot()

    if cfg.save_log:
        robot._log_manager.saveLog()

    if cfg.visualizer:
        robot.killManipulatorVisualizer()
\end{minted}

First one specifies the imports. The library forces you to type out the 
whole import path
to force you to get familiar with where different controllers or 
utilities are implemented.
This is so because you are then one step closer to 
looking at the code you are importing,
which you are highly encouraged to do.
Very often, you will get strange looking behaviour due to a 
particular controller's features or limitations.
The only way to fix those is to understand which code is actually being run,
and then either changing you approach, or writing new controllers.

Following any potential class or function definitions comes the main
function where you load the config and the robot, 
and run the desired controllers one after the other.
This architecture introduces some delay between controllers,
but it is only around $\approx$100ms.
If this ruins the robot's performance, you are almost
certainly applying an incorrect control strategy ---
the robot should be motionless when a controller exits.
Here it's very important to note that you can switch 
between different controller's within one control loop,
if control-switching behaviour is what you want
(just write a control loop that calls the other control loops 
according to desired criteria).


At the end of the script, you have to put in the following:
\begin{itemize}
		\item \fbox{\texttt{stopRobot()}} 
		\item \fbox{\texttt{log\_manager.saveLog()}} 
		\item \fbox{\texttt{killManipulatorVisualizer()}} 
\end{itemize}
You will de facto always want these functions, and they have not been automated away otherwise.
This is one of the drawbacks of keeping the core of the library as simple as possible.

\subsubsection{Dealing with parameters and settings}
Nearly all controllers come with a set of parameters which need to be
tuned for the task at hand.
Sometimes we need to run the same main file, and just change
the value of a few parameters.
This is ideally achieved through argument passing.
Depending on the situation, one needs to change many settings
and parameters (and then tune them one at a time).
To avoid having to type a very long sequence of arguments,
they should be store in a local config.yaml file.

SMC offers these capabilities in the following way.
In total, there is a dizzying array of different parameters
and settings to set. Having them all outputed every time
would make the \fbox{\texttt{--help}} option useless.
Hence, they are loaded dynamically based on what's impored
through a bit of Python magic.
With a little more magic, the dynamically loaded 
\fbox{\texttt{ConfigClass}}es also dynamically create an argument
parser.
In total, parameters are handled as follows:
\begin{enumerate}
\item Various ConfigClasses are loaded dynamically and
	put into \fbox{\texttt{GlobalConfig}}. Every parameter
	and setting is put to a sane default value.
\item Optionally, a config file is read by providing
	\fbox{\texttt{--config\_file ./path/to/file}},
	and the appropriate values are overriden.
\item Finally, the arguments are processed and take ultimate
	precedence.
\end{enumerate}

\paragraph{Defining a config class}
When developing a controller or a project,
it is also extremely likely that you will want to define 
your own ConfigClass, which will also generate appropriate arguments.
This is very simple to do:
\begin{minted}[mathescape, linenos]{python}
from smc.bookkeeping.registry import ConfigRegistry
from dataclasses import dataclass, field

ConfigRegistry.register("your_control")
@dataclass
class ConfigYour:
    arg1: str = field(default="", metadata={"help": "arg1 does ..."})
    arg2: int = field(default=0, metadata={"help": "arg2 does ..."})
\end{minted}
When running your main file, if \fbox{\texttt{ConfigYour}} is imported
at any point, the appropriate arguments will be generated.
The values are accessible from a \fbox{\texttt{GlobalConfig}} instance
cfg via the registered name, in this case
\fbox{\texttt{cfg.your\_config.arg1}}.

On this note, as a rule of thumb, nearly all magic numbers 
should be arguments. For example: 
\begin{itemize}
	\item controller gains
	\item regularization parameters, ex. Levenberg-Marquadt damping
	\item numerics parameters, ex. the minimum detectable force
	\item simulation parameters, ex. added simulated noise amplitude
	\item problem settings, ex. which interpolator you are using
	\item ...
\end{itemize}
This will allow you to quickly tune your control strategy, 
which you will practically speaking always have to do,
regardless of the particular application.



\subsection{Writing a new controller}
Copy-paste the most similar existing controller as a starting point.
Put it in the main file.

If you end up writing 3 or more controllers for the same control problem,
write a ControlLoopTemplate for the problem.
This will minimize the amount of code you have, and the you will avoid 
code copy-pasting. 
Both of these are fundamental software engineering principles,
following of which will give you the following benefits:
\begin{itemize}
\item If you need to change the structure of the problem, you
		change it in only one place. This will save some time,
		but more importantly it will \textbf{prevent hidden logic errors}
		stemming from not updating all instances 
		(this is practically guaranteed to happen otherwise).
\item It will force you to think about the structure of your problem
		and to disentangle components. This will increase your 
		understanding of the problem, and which structure is responsible for what
		(ex. what is the planner's job and what is the controllers job).
		Finally, it will make it easier to test performance.
\end{itemize}

\subsection{Assessing performance}
Quality of a robot's performance can often be visualize determined,
which is one of the reasons to use the visualizer.
To aid in this process, you should consider expanding on default visualization
which includes just the robot. Plotting paths and obstacles
will go a long way in seeing where the controller fails.
Numeric criteria start being useful only when performance differences
are not visually obvious.
More on visualization in the visualization section.

At some point, generic log analyzers will be implemented. 
At the moment, you are left to your own devices when it comes
to numeric performance measurements. 
In any event, you will want to save logs
to compare different runs of your experiment (both visually and numerically).
The starting of points are in 
\fbox{\texttt{examples/log\_analysis}}.

\subsection{Preparing for a hardware experiment}
The experiment should look same way in real life as it does
in simulation, or it will be different in an expected or reasonable way.
However, this will not be the case unless specific preparatory steps are taken.
To repeat:
\newline
\vspace{10pt}
\noindent\fbox{
\parbox{\textwidth}{
\textbf{THE EXPERIMENT WILL GO BADLY, 
POTENTIALLY RESULTING IN DAMAGES AND INJURY,
UNLESS YOU FOLLOW THE CHECKLIST BELOW! }
}}
\vspace{10pt}

\vspace{10pt}
\fbox{
\parbox{\textwidth}{
\underline{Hardware experiment preparation checklist}
\begin{enumerate}
\item Have a working controller in simulation (looks OK in the visualizer).
\item Make sure the execution is fast enough with \fbox{\texttt{--verbose}} - 
	this prints out control frequency deadline misses - a few are OK as 
	we don't provide real-time guarantees, but
	continuous printouts mean you have a delay.
\item If interacting with the environment, write functionality to simulate 
	the interaction, ex. add artificial readings to the force sensor 
	in the control loop to simulate contact,
	and make sure you get the expected behaviour.
\item Connect to the robot. Run the experiment in simulation with \fbox{\texttt{--start-from-current-pose}} argument
		to see the concrete expected motion of the robot.
\item When running the experiment:
\begin{itemize}
	\item If running for the first time, use a smaller maximum velocity with
		\fbox{\texttt{--max-v-percentage=0.1}} or similar.
	\item \textbf{HAVE YOUR HAND ABOVE THE EMERGENCY STOP BUTTON FROM EXPERIMENT START TO FINISH.}
	\item \textbf{STOP THE ROBOT UPON ANY UNEXPECTED BEHAVIOUR.} 
		You are most likely not modelling an important part of the experiment,
		or there was an issue in the communication to the robot.
\end{itemize}
\item If the experiment fails:
\begin{itemize}
	\item Inspect the logs from the simulation, and from the hardware experiment.
	If everything looks similar, you are probably not logging all relevant quantities.
	Introduce them and repeat the experiment.
\item If the problem persists, you might have a hardware fault. Run the simplest possible
		scripts until you find the same problem.
\item If you can not figure out what the problem is, contact a person responsible for hardware
		and explain the problem to them as exactly as possible.
\end{itemize}
\end{enumerate}
}}
\vspace{10pt}

